\documentclass[12pt]{article}
\usepackage{amsfonts,amsthm,amsmath,amssymb,mathtools}
\usepackage{mathrsfs}
\usepackage{tikz-cd}
\usepackage{geometry}
\usepackage{mlmodern}
\usepackage{enumitem}

\usetikzlibrary{calc,arrows.meta}

\setlist[enumerate]{noitemsep, topsep=0pt}
\setlist[itemize]{noitemsep, topsep=0pt}

\geometry{letterpaper, portrait, margin=1in}

\setcounter{section}{0}
\setcounter{secnumdepth}{1}

\renewcommand{\thesection}{\S\arabic{section}}
\renewcommand{\thesubsection}{\arabic{subsection}}
% \newcommand{\dd}{\mathop{}\!\mathrm{d}}
\newcommand{\dd}{\mathop{}\!d}

\newtheorem{thm}{Theorem}
\newenvironment{theorem}[1]{
	\renewcommand{\thethm}{#1}
	\begin{thm}
		}{\end{thm}}

\newtheorem{lma}{Lemma}
\newenvironment{lemma}[1]{
	\renewcommand{\thelma}{#1}
	\begin{lma}
		}{\end{lma}}

\theoremstyle{definition}
\newtheorem{ex}{Exercise}
\newenvironment{exercise}[1]{
	\renewcommand{\theex}{#1}
	\begin{ex}
		}{\end{ex}}

\title{Calculus III Cheat Sheet}
\date{}
\author{}

\begin{document}

\maketitle

\section{3D Space}
% \subsection{The 3D Coordinate System}

\subsection{Equations of Lines}
\begin{align*}
	\text{Vector form     :\hspace{1em}} & \vec{r} = \vec{r_0} + t\vec{v}= \left\langle x_0, y_0, z_0 \right\rangle + t \langle a, b, c \rangle \\
	\text{Parametric form :\hspace{1em}} & x = x_0 + ta, y = y_0 + tb, z = z_0 + tc                                                             \\
	\text{Symmetric form  :\hspace{1em}} & \frac{x - x_0}{a} = \frac{y - y_0}{b} = \frac{z - z_0}{c}
\end{align*}

\subsection{Equations of Planes}
Given a normal vector $\vec{n} = \langle a, b, c \rangle$ and a point on the plane $\vec{r_0}$, we have \begin{align*}
	\text{Vector equation:\hspace{1em}} & \vec{n} \cdot \left( \vec{r} - \vec{r_0} \right) = 0 \\
	\text{Scalar equation:\hspace{1em}} & ax + by + cz = d
\end{align*}

\subsection{Quadric Surfaces}
A \textit{quadric surface} is the graph of an equation in this form: \begin{equation*}
	Ax^2 + By^2 + Cz^2 + Dxy + Exz + Fyz + Gx + Hy + Iz + J = 0
\end{equation*}
Some common examples: \begin{align*}
	\text{Ellipsoid:\hspace{0.5em}}                 & \frac{x^2}{a^2} + \frac{y^2}{b^2} + \frac{z^2}{c^2} = 1  \\
	\text{Cone:\hspace{0.5em}}                      & \frac{x^2}{a^2} + \frac{y^2}{b^2} = \frac{z^2}{c^2}      \\
	\text{Cylinder:\hspace{0.5em}}                  & \frac{x^2}{a^2} + \frac{y^2}{b^2} = 1                    \\
	\text{Hyperboloid of one sheet:\hspace{0.5em}}  & \frac{x^2}{a^2} + \frac{y^2}{b^2} - \frac{z^2}{c^2} = 1  \\
	\text{Hyperboloid of two sheets:\hspace{0.5em}} & -\frac{x^2}{a^2} - \frac{y^2}{b^2} + \frac{z^2}{c^2} = 1 \\
	\text{Elliptic Paraboloid:\hspace{0.5em}}       & \frac{x^2}{a^2} + \frac{y^2}{b^2} = \frac{z}{c}          \\
	\text{Hyperbolic Paraboloid:\hspace{0.5em}}     & \frac{x^2}{a^2} - \frac{y^2}{b^2} = \frac{z}{c}          \\
\end{align*}

\subsection{Functions of Several Variables}
One way to visualize multivariable functions is to take \textit{level sets} (sometimes called \textit{level curves}, \textit{level surfaces}, \textit{contour curves}, etc. when appropriate), where you set one variable constant and see what happens with the others.
These can also be called \textit{traces} when $x$ or $y$ is what you're holding constant.

\subsection{Vector-Valued Functions}
A \textit{vector-valued function} (or \textit{vector function}) is a function which takes in one or more variables and returns a vector.
A single-variable vector-valued function is often expressed like $\vec{r}(t) = \left\langle f(t), g(t), h(t) \right\rangle$, where $f$, $g$, and $h$ are called \textit{component functions}.

\subsection{Calculus with Vector-Valued Functions}
A \textit{smooth curve} is a curve $\vec{r}'(t)$ for which $\vec{r}'(t)$ is continuous and $\vec{r}'(t)\neq 0$ for all $t$ except possibly at the endpoints.

Some useful facts (framed in $\mathbb{R}^3$, but can be generalized to $\mathbb{R}^n$): \begin{align*}
	\lim_{t \to a} \left\langle f(t), g(t), h(t) \right\rangle & = \left\langle \lim_{t \to a} f(t), \lim_{t \to a} g(t), \lim_{t \to a} h(t) \right\rangle \\
	\frac{d}{dt}\left\langle f(t), g(t), h(t) \right\rangle    & = \left\langle f'(t), g'(t), h'(t) \right\rangle                                           \\
	\int \left\langle f(t), g(t), h(t) \right\rangle dt        & = \left\langle \int f(t) dt, \int g(t) dt, \int h(t) dt \right\rangle                      \\
	\int_a^b \left\langle f(t), g(t), h(t) \right\rangle dt    & = \left\langle \int_a^b f(t) dt, \int_a^b g(t) dt, \int_a^b h(t) dt \right\rangle.
\end{align*}

\subsection{Tangent, Normal, and Binormal Vectors}
Given a function $\vec{r} : \mathbb{R} \to \mathbb{R}^3$, the \textit{unit tangent vector} is defined by \begin{equation*}
	\vec{T}(t) = \frac{\vec{r}'(t)}{\left|\left|\vec{r}'(t)\right|\right|}.
\end{equation*}
The \textit{unit normal vector} is orthogonal to the unit tangent vector and is defined by \begin{equation*}
	\vec{N}(t) = \frac{\vec{T}'(t)}{\left|\left|\vec{T}'(t)\right|\right|}.
\end{equation*}
The \textit{binormal vector} is orthogonal to both the unit tangent vector and the unit normal vector and is defined by \begin{equation*}
	\vec{B}(t) = \vec{T}(t) \times \vec{N}(t).
\end{equation*}

\subsection{Arc Length with Vector Functions}
Arc length of a single-variable, vector-valued function $\vec{r}$: \begin{equation*}
	L = \int_a^b \left| \left| r'(t) \right| \right|dt.
\end{equation*}
The \textit{arc length function} is defined as \begin{equation*}
	s(t) = \int_0^t \left| \left| r'(u) \right| \right| du,
\end{equation*}
which can be useful when we want to reparametrize the function in terms of arc length by solving the resulting $s(t)$ for $t$.

\subsection{Curvature}
The \textit{curvature} of a function is defined as \begin{equation*}
	\kappa = \left| \left| \frac{d\vec{T}}{ds} \right| \right|.
\end{equation*}
Here are two alternate forms that are often somewhat easier to compute: \begin{align*}
	\kappa & = \frac{\left| \left| \vec{T}'(t) \right| \right|}{\left| \left| r'(t) \right| \right|} & \kappa = \frac{\left| \left| r'(t)\times r''(t) \right| \right|}{\left| \left| r'(t) \right| \right|^3}
\end{align*}

\subsection{Velocity and Acceleration}
It is common to break up acceleration into its tangential and normal components ($a_T$ and $a_N$, respectively).
Letting $\vec{v}$ be the velocity and $v$ = $||\vec{v}(t)||$, these have simple formulas: \begin{align*}
	 & a_T = v' = \frac{r'(t) \cdot r''(t)}{||r'(t)||} & a_N = \kappa v^2 = \frac{||r'(t) \times r''(t)||}{||r'(t)||}
\end{align*}

\subsection{Cylindrical Coordinates}
Cylindrical coordinates are in the form $(r, \theta, z)$, where $r$ and $\theta$ act just as in polar coordinates, and the $z$ acts like the usual Cartesian $z$ (refer to Figure \ref{fig:spherical-cylindrical} below).
\subsubsection{Conversion formulas}
Cylindrical to Cartesian: \begin{align*}
	x & = r\cos\theta \\
	y & = r\sin\theta \\
	z & = z
\end{align*}
Cartesian to cylindrical: \begin{align*}
	r      & = \sqrt{x^2+y^2}                    \\
	\theta & = \tan^{-1}\left(\frac{y}{x}\right) \\
	z      & = z
\end{align*}

\subsection{Spherical Coordinates}
Spherical coordinates are of the form $(\rho, \theta, \varphi)$, where: \begin{itemize}
	\item $\rho$ is the distance from the origin to the point (so $\rho \geq 0$).
	\item $\theta$ is the same as $\theta$ from polar/cylindrical coordinates.
	\item $\varphi$ is the angle between the positive $z$-axis and the line from the origin to the point. We restrict to $0 \leq \varphi \leq \pi$.
\end{itemize}
\subsubsection{Conversion formulas}
Spherical to cylindrical: \begin{align*}
	r      & = \rho\sin\varphi \\
	\theta & = \theta          \\
	z      & = \rho\cos\varphi
\end{align*}
Spherical to Cartesian: \begin{align*}
	x & = \rho\sin\varphi\cos\theta \\
	y & = \rho\sin\varphi\sin\theta \\
	z & = \rho\cos\varphi
\end{align*}

\begin{figure}[h] \centering
	\begin{tikzpicture}[ scale=1.2, line cap=round, line join=round, >=Stealth, every node/.style={font=\small} ] \coordinate (O) at (0,0); \coordinate (X) at (-1.55,-2.05); \coordinate (Y) at (4.35,-0.15); \coordinate (Z) at (0,4.25); \coordinate (R) at (3.05,-1.32); \coordinate (P) at (3.05,2.55); \draw[->] (O) -- ($(Z)+(0,0.15)$) node[above] {$z$}; \draw[->] (O) -- ($(Y)+(0.20,-0.01)$) node[right] {$y$}; \draw[->] (O) -- ($(X)+(-0.10,-0.14)$) node[below left] {$x$}; \draw[dashed] (O) -- (R) node[pos=.55, above right] {$r$}; \draw[dashed] (R) -- (P) node[midway, right] {$z$}; \draw (O) -- (P) node[pos=.47, above left] {$\rho$}; \fill (P) circle (1.1pt); \node[anchor=south east] at ($(P)+(1.42,0.12)$) {$(x,y,z)=(\rho,\theta,\varphi)$}; \draw[->] ($(O)+(90:0.78)$) arc[start angle=90, end angle=40, radius=0.78]; \node at ($(O)+(63:1.05)$) {$\varphi$}; \draw[->] ($(O)+(-127:0.58)$) arc[start angle=-127, end angle=-23.5, radius=0.58]; \node at ($(O)+(-76:0.88)$) {$\theta$}; \coordinate (RA) at ($(R)!0.06!(P)$); \coordinate (RB) at ($(R)!0.06!(O)$); \coordinate (RC) at ($(RA)+(RB)-(R)$); \draw (RA) -- (RC) -- (RB); \end{tikzpicture}
	\caption{A comparison of spherical and cylindrical coordinates}
	\label{fig:spherical-cylindrical}
\end{figure}


\section{Partial Derivatives}

\subsection{Partial Derivatives}
To take a derivative with respect to one variable, just take the normal derivative with respect to that variable, holding all other variables constant.

For higher-order partial derivatives, note that we have $f_{xy} = f_{yx}$ at a point whenever both mixed partials are continuous on some neighborhood of that point.
This is Clairaut's theorem.

\subsection{Differentials}
The \textit{differential} or \textit{total differential} of a two-variable function $z=f(x,y)$ is defined as \begin{equation*}
	dz = df = f_x dx + f_y dy.
\end{equation*}
This extends naturally to functions of more variables.

\subsection{Chain Rule}
The chain rule for a function $z$ of $n$ variables $x_1, x_2, \dots, x_n$ which are each functions of $m$ variables $t_1, t_2, \dots, t_m$ is defined as \begin{equation*}
	\frac{\partial z}{\partial t_i} = \frac{\partial z}{\partial x_1}\frac{\partial x_1}{\partial t_i} + \frac{\partial z}{\partial x_2}\frac{\partial x_2}{\partial t_i} + \cdots + \frac{\partial z}{\partial x_n}\frac{\partial x_n}{\partial t_i}.
\end{equation*}

\subsection{Directional Derivatives}
We define the \textit{gradient} of a three-variable function $f$ to be \begin{equation*}
	\nabla f = \left\langle f_x, f_y, f_z \right\rangle.
\end{equation*}
Now given any unit vector $\vec{u}$, we say that the \textit{directional derivative} of $f$ in the direction of $\vec{u}$ is \begin{equation*}
	D_{\vec{u}}f = \nabla f \cdot \vec{u}.
\end{equation*}
A couple of theorems about gradients: \begin{enumerate}
	\item The maximum value of $D_{\vec{u}}f(\vec{x})$ is given by $||\nabla f(\vec{x})||$ and will occur in the direction of $\nabla f(\vec{x})$.
	\item The gradient vector $\nabla f(x_0, y_0)$ is orthogonal to the level curve $f(x, y) = k$ at the point $(x_0, y_0)$.
\end{enumerate}

\section{Applications of Partial Derivatives}

\subsection{Tangent Planes}
The equation of the tangent plane to a graph $f(x, y, z) = k$ at the point $(x_0, y_0, z_0)$ is given by \begin{equation*}
	f_x(x_0, y_0, z_0)(x-x_0) + f_y(x_0, y_0, z_0)(y-y_0) + f_z(x_0, y_0, z_0)(z-z_0)  = 0,
\end{equation*}
which you can remember because it's ensuring that points on the plane, which are of the form $(x-x_0, y-y_0, z-z_0)$, are orthogonal to the gradient, so the equation becomes \begin{equation*}
	(\nabla f)(x_0,y_0,z_0) \cdot (x-x_0, y-y_0, z-z_0) = 0.
\end{equation*}

\subsection{Relative Minimums and Maximums}
The point $(a, b)$ is a \textit{critical point} of $f(x, y)$ provided one of the following is true: \begin{enumerate}
	\item $\nabla f(a,b) =\vec{0}$
	\item $f_x(a,b)$ and/or $f_y(a, b)$ doesn't exist.
\end{enumerate}
If $(a, b)$ is a relative extrema of $f(x, y)$ and the first order derivatives of $f$ exist at $(a, b)$, then $(a, b)$ is a critical point and $\nabla f (a, b) = \vec{0}$.

Let $(a, b)$ be a critical point of $f(x, y)$ where the second order partials are continuous around $(a, b)$.
Define \begin{equation*}
	D = D(a, b) = f_{xx}(a, b)f_{yy}(a, b) - \left[f_{xy}(a,b)\right]^2.
\end{equation*}
Then \begin{enumerate}
	\item If $D > 0$, and $f_{xx}(a,b)>0$ then there is a relative minimum at $(a, b)$
	\item If $D > 0$, and $f_{xx}(a,b)<0$ then there is a relative maximum at $(a, b)$
	\item If $D<0$, then the point $(a, b)$ is a saddle point
	\item If $D = 0$, then $(a, b)$ may be any of the three types.
\end{enumerate}

\subsection{Absolute Minimums and Maximums}
Extreme Value Theorem: if $f(x, y)$ is continuous on a closed and bounded set $D$, then it achieves an absolute maximum and an absolute minimum on $D$.

Method for finding absolute extrema on a closed region $D$: \begin{enumerate}
	\item Find all the critical points.
	\item Find extrema along the boundary (reduces to a Calc I problem).
	\item The largest is the absolute maximum and the smallest is the absolute maximum.
\end{enumerate}

\subsection{Lagrange Multipliers}
This is another way to optimize, called the Method of Lagrange Multipliers: \begin{enumerate}
	\item Solve the following system of equations: \begin{align*}
		      \nabla f(x, y, z) & = \lambda\nabla g(x, y, z) \\
		      g(x, y, z)        & = k.
	      \end{align*}
	\item Plug all solutions from the first step into $f$ and identify the minimum and maximum values, provided they exist and $\nabla g \neq \vec{0}$ at the point.
	      (The constant $\lambda$ is the \textit{Lagrange Multiplier}.)
	\item Look for any remaining solutions along the boundary.
\end{enumerate}

\section{Multiple Integrals}

\subsection{Double Integrals}
The definition of a double integral over a rectangular region $R$: \begin{equation*}
	\iint\limits_R f(x, y) \dd A = \lim_{n,m\to\infty} \sum_{i=1}^n\sum_{j=1}^m f(x_i^*, y_j^*) \Delta A,
\end{equation*}
where $x_i^*$ and $y_j^*$ are points from somewhere within the subdivisions (can be just midpoints).

\subsection{Iterated Integrals}
Fubini's Theorem: If $f(x, y)$ is continuous on $R = [a, b] \times [c, d]$, then \begin{equation*}
  \iint\limits_R f(x,y) \dd A = \int_a^b\int_c^d f(x,y) \dd y \dd x = \int_c^d\int_a^b f(x, y) \dd x \dd y.
\end{equation*}
Helpful trick: if $f(x, y) = g(x)h(y)$ and we are integrating over $R = [a,b] \times [c, d]$, then \begin{equation*}
  \iint\limits_R f(x,y) \dd A = \left( \int_a^b g(x) \dd x \right) \left( \int_c^d h(y) dy \right).
\end{equation*}
Note that if we take the indefinite integral of a multivariable function with respect to one variable, the ``constant'' of integration is a function of the other variables.

\subsection{Double Integrals over General Regions}

\subsection{Double Integrals in Polar Coordinates}

\subsection{Triple Integrals}

\subsection{Triple Integrals in Cylindrical Coordinates}

\subsection{Triple Integrals in Spherical Coordinates}

\subsection{Change of Variables}

\subsection{Surface Area}

\subsection{Area and Volume Revisited}

\section{Line Integrals}
.

\subsection{Vector Fields}

\subsection{Line Integrals --- Part I}

\subsection{Line Integrals --- Part II}

\subsection{Line Integrals of Vector Fields}

\subsection{Fundamental Theorem for Line Integrals}

\subsection{Conservative Vector Fields}

\subsection{Green's Theorem}

\section{Surface Integrals}
.

\subsection{Curl and Divergence}

\subsection{Parametric Surfaces}

\subsection{Surface Integrals}

\subsection{Surface Integrals of Vector Fields}

\subsection{Stokes' Theorem}

\subsection{Divergence Theorem}

\end{document}
